% GENERATED FILE. Edit canonical lessons, then run npm run generate:books.
% canonical-source-hash: fnv1a64:6f430e79089f6a30
% canonical-lessons: PT-C440-morada, PT-C440-periferia, PT-C440-freguesia, PT-C440-concelho, PT-C440-intercomunicador, PT-R440-revisao-433-440

\chapter{A morada}
\label{ch:pt-a-morada}
\hlchaptermodality{440}

\begin{chapteropening}
\textbf{By the end of this chapter:} \emph{I can find an address: the address, the outskirts, the parish, the municipality and the entry phone.}

Complete the last lesson of chapter 440: I can find an address: the address, the outskirts, the parish, the municipality and the entry phone.
\end{chapteropening}

\section[a morada]{a morada — the address (where someone lives)}

\label{lesson:PT-C440-morada}



\begin{quote}
Before the new one: say the Portuguese for the south-west, then the Portuguese for the landmark.
\end{quote}

\subsection*{You'll want to know: a morada}
\textbf{a morada} — \textquotedblleft{}the address (where someone lives)\textquotedblright{}.

\begin{quote}
Escrevi a morada no envelope.
\end{quote}

\begin{quote}
\emph{I wrote the address on the envelope.}
\end{quote}

In Brazil: \emph{o endereço}, which Portugal uses for e-mail: \emph{o endereço eletrónico}.

\begin{cousinweb}
\emph{morada} is built on \emph{morar} 'to live, to dwell', which is inherited from Latin \emph{morārī} 'to linger, to stay' (Lewis \& Short).
\end{cousinweb}

\begin{grammarlens}[title={What you've built}]
One more word for telling, explaining and describing at B1.
\end{grammarlens}

\subsection*{Your turn}
\begin{itemize}
\raggedright
  \item \emph{Say it:} \emph{a morada}
  \item \emph{Say it:} \emph{a morada}, once more
  \item \emph{Say it:} a sentence of your own with \emph{a morada}
  \item \emph{Recall:} say the Portuguese for the south-west, then the Portuguese for the landmark, then say \emph{a morada} again
\end{itemize}

\begin{tcolorbox}[breakable,colback=teal!4,colframe=teal!35!black,title={Before you move on}]
What does a morada mean? (\textquotedblleft{}The address (where someone lives)\textquotedblright{}.) Say it once more.
\end{tcolorbox}
\section[a periferia]{a periferia — the outskirts, the edge}

\label{lesson:PT-C440-periferia}



\begin{quote}
Before the new one: say the Portuguese for the landmark, then the Portuguese for the address.
\end{quote}

\subsection*{You'll want to know: a periferia}
\textbf{a periferia} — \textquotedblleft{}the outskirts, the edge\textquotedblright{}.

\begin{quote}
Muitas famílias vivem na periferia da cidade.
\end{quote}

\begin{quote}
\emph{Many families live on the outskirts of the city.}
\end{quote}

Also the edge of a circle.

\begin{cousinweb}
\emph{periferia} is a learned borrowing, through Late Latin, from Greek \emph{periphéreia} 'circumference', from \emph{perí} 'around' and \emph{phérein} 'to carry' (Lewis \& Short).
\end{cousinweb}

\begin{grammarlens}[title={What you've built}]
One more word for telling, explaining and describing at B1.
\end{grammarlens}

\subsection*{Your turn}
\begin{itemize}
\raggedright
  \item \emph{Say it:} \emph{a periferia}
  \item \emph{Say it:} \emph{a periferia}, once more
  \item \emph{Say it:} a sentence of your own with \emph{a periferia}
  \item \emph{Recall:} say the Portuguese for the landmark, then the Portuguese for the address, then say \emph{a periferia} again
\end{itemize}

\begin{tcolorbox}[breakable,colback=teal!4,colframe=teal!35!black,title={Before you move on}]
What does a periferia mean? (\textquotedblleft{}The outskirts, the edge\textquotedblright{}.) Say it once more.
\end{tcolorbox}
\section[a freguesia]{a freguesia — the parish (local district)}

\label{lesson:PT-C440-freguesia}



\begin{quote}
Before the new one: say the Portuguese for the address, then the Portuguese for the outskirts.
\end{quote}

\subsection*{You'll want to know: a freguesia}
\textbf{a freguesia} — \textquotedblleft{}the parish (local district)\textquotedblright{}.

\begin{quote}
A minha avó mora numa freguesia pequena.
\end{quote}

\begin{quote}
\emph{My grandmother lives in a small parish.}
\end{quote}

The smallest unit of local government in Portugal, run by \emph{a junta de freguesia}.

\begin{cousinweb}
\emph{freguesia} is built on \emph{freguês} 'parishioner, customer' with \emph{-ia} (EtymDB); \emph{freguês} goes back to Latin \emph{fīlius ecclēsiae} 'son of the church'.
\end{cousinweb}

\begin{grammarlens}[title={What you've built}]
One more word for telling, explaining and describing at B1.
\end{grammarlens}

\subsection*{Your turn}
\begin{itemize}
\raggedright
  \item \emph{Say it:} \emph{a freguesia}
  \item \emph{Say it:} \emph{a freguesia}, once more
  \item \emph{Say it:} a sentence of your own with \emph{a freguesia}
  \item \emph{Recall:} say the Portuguese for the address, then the Portuguese for the outskirts, then say \emph{a freguesia} again
\end{itemize}

\begin{tcolorbox}[breakable,colback=teal!4,colframe=teal!35!black,title={Before you move on}]
What does a freguesia mean? (\textquotedblleft{}The parish (local district)\textquotedblright{}.) Say it once more.
\end{tcolorbox}
\section[o concelho]{o concelho — the municipality, the council area}

\label{lesson:PT-C440-concelho}



\begin{quote}
Before the new one: say the Portuguese for the outskirts, then the Portuguese for the parish.
\end{quote}

\subsection*{You'll want to know: o concelho}
\textbf{o concelho} — \textquotedblleft{}the municipality, the council area\textquotedblright{}.

\begin{quote}
O concelho tem dez freguesias.
\end{quote}

\begin{quote}
\emph{The municipality has ten parishes.}
\end{quote}

Not \emph{conselho}, advice, which sounds the same.

\begin{cousinweb}
\emph{concelho} is inherited from Latin \emph{concilium} 'council, meeting' (EtymDB).
\end{cousinweb}

\begin{grammarlens}[title={What you've built}]
One more word for telling, explaining and describing at B1.
\end{grammarlens}

\subsection*{Your turn}
\begin{itemize}
\raggedright
  \item \emph{Say it:} \emph{o concelho}
  \item \emph{Say it:} \emph{o concelho}, once more
  \item \emph{Say it:} a sentence of your own with \emph{o concelho}
  \item \emph{Recall:} say the Portuguese for the outskirts, then the Portuguese for the parish, then say \emph{o concelho} again
\end{itemize}

\begin{tcolorbox}[breakable,colback=teal!4,colframe=teal!35!black,title={Before you move on}]
What does o concelho mean? (\textquotedblleft{}The municipality, the council area\textquotedblright{}.) Say it once more.
\end{tcolorbox}
\section[o intercomunicador]{o intercomunicador — the entry phone, the intercom}

\label{lesson:PT-C440-intercomunicador}



\begin{quote}
Before the new one: say the Portuguese for the parish, then the Portuguese for the municipality.
\end{quote}

\subsection*{You'll want to know: o intercomunicador}
\textbf{o intercomunicador} — \textquotedblleft{}the entry phone, the intercom\textquotedblright{}.

\begin{quote}
Toquei no intercomunicador, mas ninguém respondeu.
\end{quote}

\begin{quote}
\emph{I rang the entry phone, but nobody answered.}
\end{quote}

Many people simply say \emph{a campainha}.

\begin{cousinweb}
\emph{intercomunicador} is built from \emph{inter-} 'between' and \emph{comunicar}, a learned borrowing from Latin \emph{commūnicāre} 'to share' (Lewis \& Short).
\end{cousinweb}

\begin{grammarlens}[title={What you've built}]
One more word for telling, explaining and describing at B1.
\end{grammarlens}

\subsection*{Your turn}
\begin{itemize}
\raggedright
  \item \emph{Say it:} \emph{o intercomunicador}
  \item \emph{Say it:} \emph{o intercomunicador}, once more
  \item \emph{Say it:} a sentence of your own with \emph{o intercomunicador}
  \item \emph{Recall:} say the Portuguese for the parish, then the Portuguese for the municipality, then say \emph{o intercomunicador} again
\end{itemize}

\begin{tcolorbox}[breakable,colback=teal!4,colframe=teal!35!black,title={Before you move on}]
What does o intercomunicador mean? (\textquotedblleft{}The entry phone, the intercom\textquotedblright{}.) Say it once more.
\end{tcolorbox}
\section[(dialogue)]{Revisão dos capítulos 433–440}

\label{lesson:PT-R440-revisao-433-440}



\begin{quote}
Say the Portuguese for the racket and to creak.
\end{quote}

\subsection*{Your turn}
Say these aloud:

\begin{itemize}
\raggedright
  \item the racket, to creak, the cockroach, mould, dirt
  \item to complain, the complaint, to grumble, to regret, the compensation
  \item to guide, the compass, lost, the travel plan, to locate
  \item the directions, the roundabout, to go round, the road junction, to carry on
  \item the avenue, the side street, the alley, the narrow lane, the fork
  \item the climb, the way down, the steep street, to run alongside, to come out
  \item the north-east, the north-west, the south-east, the south-west, the landmark
  \item the address, the outskirts, the parish, the municipality, the entry phone
\end{itemize}

\begin{tcolorbox}[breakable,colback=teal!4,colframe=teal!35!black,title={Before you move on}]
The racket? (\textbf{a algazarra}.) To creak? (\textbf{ranger}.) The cockroach? (\textbf{a barata}.) Mould? (\textbf{o bolor}.) Dirt? (\textbf{a sujidade}.) To complain? (\textbf{reclamar}.) The complaint? (\textbf{a reclamação}.) To grumble? (\textbf{resmungar}.) To regret? (\textbf{lamentar}.) The compensation? (\textbf{a indemnização}.) To guide? (\textbf{orientar}.) The compass? (\textbf{a bússola}.) Lost? (\textbf{desorientado}.) The travel plan? (\textbf{o roteiro}.) To locate? (\textbf{localizar}.) The directions? (\textbf{as indicações}.) The roundabout? (\textbf{a rotunda}.) To go round? (\textbf{contornar}.) The road junction? (\textbf{o entroncamento}.) To carry on? (\textbf{prosseguir}.) The avenue? (\textbf{a avenida}.) The side street? (\textbf{a travessa}.) The alley? (\textbf{o beco}.) The narrow lane? (\textbf{a viela}.) The fork? (\textbf{a bifurcação}.) The climb? (\textbf{a subida}.) The way down? (\textbf{a descida}.) The steep street? (\textbf{a ladeira}.) To run alongside? (\textbf{ladear}.) To come out? (\textbf{desembocar}.) The north-east? (\textbf{o nordeste}.) The north-west? (\textbf{o noroeste}.) The south-east? (\textbf{o sudeste}.) The south-west? (\textbf{o sudoeste}.) The landmark? (\textbf{o marco}.) The address? (\textbf{a morada}.) The outskirts? (\textbf{a periferia}.) The parish? (\textbf{a freguesia}.) The municipality? (\textbf{o concelho}.) The entry phone? (\textbf{o intercomunicador}.)
\end{tcolorbox}
